\chapter{Вещественно-скалярное скручивание: представление KO6 и проверка знака}

Предыдущая проверка выбрала обмен двух вещественных скалярных копий. Здесь
алгебра представляется на существующем 18-мерном гильбертовом пространстве
KO6, условия скрученного первого порядка проверяются на полном базисе, а
опора флуктуаций вычисляется до анализа квартичного знака.

\section{Скрученная алгебра и её представление}

Пусть
\begin{equation}
 \mathcal A_{\rm tw}^{(0)}=\mathbb R_{0,+}\oplus\mathbb R_{0,-}
 \oplus M_3(\mathbb R)_G\oplus\mathbb C_2,
\end{equation}
а инволютивный автоморфизм равен
\begin{equation}
 \rho(\lambda_+,\lambda_-,A,z)=(\lambda_-,\lambda_+,A,z).
 \label{eq:v5-real-flip-rho}
\end{equation}
При порядке блоков \(p_0,p_1,p_2,c_0,c_1,c_2\) левое представление имеет вид
\begin{equation}
 \pi_{\rm tw}(\lambda_+,\lambda_-,A,z)=\operatorname{diag}
 (\lambda_+I_3,A,A,\lambda_-I_3,\lambda_+I_3,\bar zI_3).
 \label{eq:v5-real-flip-representation}
\end{equation}
Оно задаётся исходными метками и градуировкой, имеет вещественный линейный
ранг \(1+1+9+2=13\) и является точным. Вещественная структура определяет
противоположное действие обычным способом.

\section{Вещественные скрученные условия порядка}

Для вещественной скрученной тройки
\begin{equation}
 \rho^\circ(b^\circ)=(\rho^{-1}(b))^\circ,
\end{equation}
а условие первого порядка равно
\begin{equation}
 [[D,\pi(a)]_\rho,J\pi(b^*)J^{-1}]_{\rho^\circ}=0.
 \label{eq:v5-real-twisted-first-order}
\end{equation}
Поскольку обмен инволютивен, это условие разворачивается в
\begin{multline}
 (D\pi(a)-\pi(\rho(a))D)J\pi(b^*)J^{-1}\\
 -J\pi(\rho(b^*))J^{-1}(D\pi(a)-\pi(\rho(a))D)=0.
 \label{eq:v5-real-twisted-first-order-expanded}
\end{multline}
Это формула (2.12) из \path{arXiv:1601.00219}.

Для полного вещественного базиса из тринадцати элементов оба остатка равны
нулю:
\begin{align}
 \max_{a,b}\|[\pi(a),J\pi(b^*)J^{-1}]\|_F&=0,\\
 \max_{a,b}\|[[D,\pi(a)]_\rho,J\pi(b^*)J^{-1}]_{\rho^\circ}\|_F&=0.
\end{align}
Без скручивания условие первого порядка для удвоенной алгебры нарушается.

\begin{theorem}[Существование скрученного вещественно-скалярного представления]
Избирательное удвоение вещественного скаляра имеет точное представление на
существующем пространстве KO6 и удовлетворяет условиям нулевого и
скрученного первого порядка без добавления фермионов.
\end{theorem}

\section{Точная опора скрученных 1-форм}

\begin{equation}
 \Omega^1_{D,\rho}=\operatorname{span}
 \{\pi(a)[D,\pi(b)]_\rho:a,b\in\mathcal A_{\rm tw}^{(0)}\}.
\end{equation}
Комплексный ранг этого пространства равен двадцати, а блочная опора точно
равна
\begin{equation}
 (p_0,p_1),(p_1,p_0),(c_1,c_2),(c_2,c_1).
 \label{eq:v5-real-flip-oneform-support}
\end{equation}
Флуктуированный оператор \(D_{A_\rho}=D+A_\rho+JA_\rho J^{-1}\) сохраняет
исходные ближайшие рёбра и не создаёт диагонального блока или блока между
концами пути. На частичной половине он имеет вид
\begin{equation}
 D_{A_\rho,p}=\begin{pmatrix}
 0&X_\rho^\dagger&0\\X_\rho&0&Y_\rho^\dagger\\0&Y_\rho&0
 \end{pmatrix}.
 \label{eq:v5-real-flip-fluctuated-chain}
\end{equation}

\section{Сохранение препятствия по знаку}

Положим \(A_\rho=X_\rho X_\rho^\dagger\),
\(B_\rho=Y_\rho^\dagger Y_\rho\). Для любого такого самосопряжённого
оператора
\begin{equation}
 \Tr D_{A_\rho,p}^4=2\Tr(A_\rho+B_\rho)^2.
 \label{eq:v5-real-flip-positive-cross}
\end{equation}
Смешанный член положителен, тогда как требуется
\(\Tr(A_\rho-B_\rho)^2\). На радиальном сечении вновь сравниваются
\begin{equation}6(\rho^2+r^2)^2\qquad\hbox{и}\qquad(\rho^2-r^2)^2.\end{equation}
Следовательно, скрученная геометрия существует, но обычный спектральный след
флуктуированного оператора не порождает потенциал отображения момента.

\begin{nogocriterion}[Недостаточность скручивания без нового следа]
Изменение дифференциального исчисления при сохранении обычного невзвешенного
\(\Tr f(D_{A_\rho}^2)\) не восстанавливает ориентационные сведения,
утраченные спектром того же трёхвершинного графа. Требуется отдельно
выведенный \(\rho\)-след, модулярный вес или иной функционал кривизны.
\end{nogocriterion}

\section{Вывод}

\begin{protocol}
\begin{enumerate}
 \item точное представление четырёх слагаемых: \textbf{получено};
 \item условия нулевого и скрученного первого порядка KO6:
 \textbf{выполнены};
 \item новая диагональная или концевая кривизна: \textbf{отсутствует};
 \item требуемый смешанный знак: \textbf{не получен};
 \item вещественно-скалярный обмен с обычным следом:
 \textbf{динамически закрыт};
 \item физическое замыкание: \textbf{не достигнуто}.
\end{enumerate}
\end{protocol}

Следующая проверка --- \path{version5_flip_twisted_trace_positivity_gate}.
Машинный сертификат сохранён в
\path{s2t_v5_real_scalar_flip_twisted_ko6_gate_results.json}.